% Section 6: Claim boundaries and the remaining programme.

\subsection{The claim ledger}

The following table is the paper's ledger.  It records the status of every substantive
statement; the governing file is
\path{governance/00e-paper-VI-process-geometry-scope.md}, and the status labels are
those of \path{governance/05c-paper-VI-status-register.md}.

\begin{center}
\footnotesize
\begin{longtable}{@{}>{\raggedright\arraybackslash}p{0.42\linewidth}>{\raggedright\arraybackslash}p{0.22\linewidth}>{\raggedright\arraybackslash}p{0.30\linewidth}@{}}
\toprule
Statement & Status & Evidence \\
\midrule
open finite program, construction order, schedule
  & definition & process corpus \cite{AdvaProcessCore2026,ProcessGeometryCore2026} \\
three readings $K,C,F$ of one process
  & \texttt{CLASSICAL} (import) & process corpus \\
slice composition and schedule independence
  & \texttt{CLASSICAL} (import) & process corpus, finite fixtures \\
arity-one principle (slot order from the carrier)
  & \texttt{STRUCTURAL PROPOSAL} & this paper, \S\ref{sec:processes} \\
sequential histories are single-spine programs
  & \texttt{PROVED} & Proposition~\ref{prop:sequential-leaf} \\
canonical evaluation order, all plane trees
  & \texttt{PROVED} & Theorem~\ref{thm:canonical-order} \\
canonical order, $3\le n\le 6$ (193 trees)
  & \texttt{COMPUTATIONALLY VERIFIED EXAMPLE} & the canonical-choice verifier (\S\ref{sec:orders}) \\
screened alignment: carousel is not $J$
  & \texttt{PROVED} (correspondence record) & \S\ref{sec:carousel} \\
Laplacian does not carry the index
  & \texttt{PROVED} (finite obstruction) & No-go~\ref{nogo:laplacian-index} \\
screw module $C^{4}=S$
  & \texttt{COMPUTATIONALLY VERIFIED EXAMPLE} & \texttt{verify\_planar.py} \\
carousel as \emph{the} carrier of the screw
  & \texttt{STRUCTURAL PROPOSAL} & \S\ref{sec:carousel} \\
sequential recovery equivalence
  & \texttt{PROVED} & Theorem~\ref{thm:sequential-recovery} \\
general process--expression equivalence
  & \texttt{CONJECTURE} & Conjecture~\ref{conj:general-equivalence} \\
one-hole contexts as compiled observations
  & \texttt{OPEN} & \S\ref{sec:equivalence} \\
charge as frontier data
  & \texttt{OPEN} & \S\ref{sec:equivalence} \\
\bottomrule
\end{longtable}
\end{center}

\subsection{Open problems}

\begin{problem}[The general functor]\label{prob:general-functor}
Construct the functor of Conjecture~\ref{conj:general-equivalence} and prove that it
preserves the canonical order, or refute the conjecture by a tree whose schedule
structure has no process image.
\end{problem}

\begin{problem}[Derive the shift]\label{prob:derive-shift}
Derive the index shift $S$ from the arithmetic or process data rather than postulating
it.  The candidate carriers supplied by the cut Laplacian are excluded by
No-go~\ref{nogo:laplacian-index}; the shift must come from a planar datum.
\end{problem}

\begin{problem}[Absorb the base corner]\label{prob:absorb-sr}
Extend the planar module so that the base corner $SR$ is represented rather than
added as an entry, and check whether the extended module still satisfies
$C^{4}=S$.
\end{problem}

\begin{problem}[Canonicality of the criterion]\label{prob:criterion}
Decide whether another natural criterion --- for example preserving the left-to-right
order of leaves, or invariance under the mirror --- selects the same unique linear
extension as Theorem~\ref{thm:canonical-order}.  Agreement would strengthen the
canonicality; disagreement would show that the choice depends on the criterion and
must be declared.
\end{problem}

\begin{problem}[Two layers, one transform]\label{prob:layers}
State precisely the relation between the planar canonicalisation of
Theorem~\ref{thm:canonical-order} and the poset-layer obstruction for the transform
$\tau$ (\texttt{OQ-098}): the first chooses one linear extension, the second says that
$\tau$ does not descend to the down-set lattice.  The two statements are different and
the difference must be written down, not conflated.
\end{problem}

\begin{problem}[Fringing and orbit growth]\label{prob:fringing}
Decide whether the carousel effect, the fringing law $\Delta(Q(T))=3\Delta(T)$ and the
orbit growth $O(n^{\log_{2}3})$ correspond to arithmetic periodic orbits, as the
research notes propose.  This needs the arithmetic orbit structure and is not settled
by the carrier correspondence.
\end{problem}

\begin{problem}[Program-side fixtures]\label{prob:fixtures}
Store the three acceptance tests of \S\ref{sec:equivalence} as executed fixtures on
the program side: tree census, schedule count, and carrier check.  Until they exist,
the equivalence remains conjectural even where each side has been checked separately.
\end{problem}

\begin{problem}[Presentation and objectification]\label{prob:presentation}
Bring the process corpus's presentation and objectification layers into the paper's
scope, or state explicitly that they are outside it.  The present draft imports only
the native process, the three readings, and the slice laws; the rest of the Process
Geometry programme is cited, not developed.
\end{problem}

\subsection{What this paper does not claim}

The paper does not claim that every program is an arithmetic expression, that values
determine programs, that the equivalence is proved in general, that the carousel
correspondence is canonical, that the screw module is the unique carrier, or that any
complexity or decidability consequence follows from the correspondence.  None of the
proved statements of Papers~0--V is redefined, strengthened, or weakened here, and no
proved node of those papers may depend on this one.

\subsection{The remaining programme}

The programme has a clear order of dependence.  First the general functor
(Problem~\ref{prob:general-functor}), because everything else is stated through it.
Then the derivation of the shift (Problem~\ref{prob:derive-shift}) and the absorption
of the base corner (Problem~\ref{prob:absorb-sr}), which decide whether the planar
module is the right carrier.  Then the criterion question
(Problem~\ref{prob:criterion}), which decides whether canonicality survives a change
of presentation.  The arithmetic directions --- fringing and orbits
(Problem~\ref{prob:fringing}) --- come last, because they presuppose the arithmetic
orbit structure and are not settled by the carrier alone.

The process corpus is cited at a declared revision.  A later revision that changes the
slice laws or the observer policy is a boundary of this paper, not an automatic
extension of it: the paper states what it imported and when.
